\documentclass[10pt]{extarticle}
\usepackage[a4paper,margin=1.5cm]{geometry}
\usepackage[T1]{fontenc}
\usepackage{amsmath,amssymb}
\usepackage{graphicx}
\usepackage[export]{adjustbox}
\usepackage{multicol}
\usepackage{enumitem}
\usepackage{booktabs,array,tabularx}
\setlength{\columnsep}{0.7cm}
\setlength{\parindent}{0pt}
\setlength{\parskip}{3pt}
\pagestyle{empty}
\makeatletter
\newcommand{\qfigdef}[4]{\@namedef{qf@file@#1}{#2}\@namedef{qf@w@#1}{#3}\@namedef{qf@wide@#1}{#4}}
\newcommand{\qf@check}[1]{\@ifundefined{qf@file@#1}{\PackageError{score2}{Unknown figure id #1}{}}{}}
\newcommand{\qf@img}[4][]{\includegraphics[width=#2\dimexpr\@nameuse{qf@w@#4}pt\relax,max width=#3,max height=0.45\textheight,#1]{\@nameuse{qf@file@#4}}}
\newif\ifqf@full
\newif\ifqf@boxed
\newsavebox{\qf@box}
\newcommand{\qf@open}{\global\qf@fulltrue\par\global\qf@boxedtrue
  \setbox\qf@box\vbox\bgroup\hsize\textwidth\linewidth\textwidth}
\newcommand{\qf@emit}[1]{\par\egroup\global\qf@boxedfalse
  \ifdim\dimexpr\ht\qf@box+\dp\qf@box\relax<#1\textheight\noindent\usebox{\qf@box}\par\else\unvbox\qf@box\fi}
\newcommand{\qfullsplit}{\ifqf@boxed\qf@emit{0.85}\fi}
\newcommand{\qf@close}{\ifqf@boxed\qf@emit{0.45}\fi\global\qf@fullfalse}
\newcommand{\qfullbegin}{\end{multicols}\qf@open}
\newcommand{\qfullend}{\qf@close\begin{multicols}{2}}
\newcommand{\qfullstart}{\qf@open}
\newcommand{\qfulldone}{\qf@close}
\newcommand{\qfig}[2][1]{\qf@check{#2}%
  \ifqf@full\qf@blk{#1}{#2}\else
  \ifnum\@nameuse{qf@wide@#2}=1 \expandafter\@firstoftwo\else\expandafter\@secondoftwo\fi
  {\end{multicols}\begin{center}\qf@img{#1}{\textwidth}{#2}\end{center}\begin{multicols}{2}}%
  {\qf@blk{#1}{#2}}\fi}
\newcommand{\qf@blk}[2]{\par\smallskip{\centering\qf@img{#1}{\linewidth}{#2}\par}\smallskip}
\newcommand{\qfigi}[2][1]{\qf@check{#2}\qf@img[valign=c]{#1}{\linewidth}{#2}}
\makeatother
\newcommand{\papersec}[1]{\par\bigskip{\centering\Large\bfseries PAPER-#1\par}\medskip}
\newcommand{\partsec}[2]{\par\medskip\noindent{\large\bfseries PART-#1: #2}\par\smallskip}
\newcommand{\sect}[1]{\par\smallskip\noindent\textbf{#1}\par}
\newcommand{\qn}[1]{\par\smallskip\noindent\textbf{#1.}~}
\newcommand{\paraq}[2]{\par\smallskip\noindent\textbf{Paragraph for Questions #1 and #2}\par}
\newenvironment{qoptions}{\begin{itemize}[leftmargin=1.8em,itemsep=1pt,topsep=2pt,label={}]}{\end{itemize}}
\newcommand{\opt}[2]{\item[(#1)] #2}
% figure registry (generated by score2_tools.py finalize)
\qfigdef{p1_phy_s1_q01_1}{JEE Advanced (2021 Series) - Test 12 (25 May 2021)-img/p1_phy_s1_q01_1.png}{101.5}{0}
\qfigdef{p1_phy_s1_q02_1}{JEE Advanced (2021 Series) - Test 12 (25 May 2021)-img/p1_phy_s1_q02_1.png}{91.4}{0}
\qfigdef{p1_phy_s1_q03_1}{JEE Advanced (2021 Series) - Test 12 (25 May 2021)-img/p1_phy_s1_q03_1.png}{102.4}{0}
\qfigdef{p1_phy_s1_q04_1}{JEE Advanced (2021 Series) - Test 12 (25 May 2021)-img/p1_phy_s1_q04_1.png}{101.5}{0}
\qfigdef{p1_phy_s1_q04_2}{JEE Advanced (2021 Series) - Test 12 (25 May 2021)-img/p1_phy_s1_q04_2.png}{101.1}{0}
\qfigdef{p1_phy_s1_q05_1}{JEE Advanced (2021 Series) - Test 12 (25 May 2021)-img/p1_phy_s1_q05_1.png}{102.3}{0}
\qfigdef{p1_phy_s1_q06_1}{JEE Advanced (2021 Series) - Test 12 (25 May 2021)-img/p1_phy_s1_q06_1.png}{208.6}{0}
\qfigdef{p1_phy_s1_q08_1}{JEE Advanced (2021 Series) - Test 12 (25 May 2021)-img/p1_phy_s1_q08_1.png}{102.6}{0}
\qfigdef{p1_phy_s1_q09_1}{JEE Advanced (2021 Series) - Test 12 (25 May 2021)-img/p1_phy_s1_q09_1.png}{128.6}{0}
\qfigdef{p1_phy_s1_q10_1}{JEE Advanced (2021 Series) - Test 12 (25 May 2021)-img/p1_phy_s1_q10_1.png}{171.3}{0}
\qfigdef{p1_phy_s2_q01_1}{JEE Advanced (2021 Series) - Test 12 (25 May 2021)-img/p1_phy_s2_q01_1.png}{104.8}{0}
\qfigdef{p1_phy_s2_q02_1}{JEE Advanced (2021 Series) - Test 12 (25 May 2021)-img/p1_phy_s2_q02_1.png}{140.6}{0}
\qfigdef{p1_phy_s2_q03_1}{JEE Advanced (2021 Series) - Test 12 (25 May 2021)-img/p1_phy_s2_q03_1.png}{104.4}{0}
\qfigdef{p1_phy_s2_q05_1}{JEE Advanced (2021 Series) - Test 12 (25 May 2021)-img/p1_phy_s2_q05_1.png}{157.9}{0}
\qfigdef{p1_phy_s2_q06_1}{JEE Advanced (2021 Series) - Test 12 (25 May 2021)-img/p1_phy_s2_q06_1.png}{135.2}{0}
\qfigdef{p1_phy_s2_q08_1}{JEE Advanced (2021 Series) - Test 12 (25 May 2021)-img/p1_phy_s2_q08_1.png}{128.0}{0}
\qfigdef{p1_phy_s2_q09_1}{JEE Advanced (2021 Series) - Test 12 (25 May 2021)-img/p1_phy_s2_q09_1.png}{102.3}{0}
\qfigdef{p1_phy_s2_q10_1}{JEE Advanced (2021 Series) - Test 12 (25 May 2021)-img/p1_phy_s2_q10_1.png}{106.5}{0}
\qfigdef{p1_chem_s1_q01_2}{JEE Advanced (2021 Series) - Test 12 (25 May 2021)-img/p1_chem_s1_q01_2.png}{24.0}{0}
\qfigdef{p1_chem_s1_q02_A}{JEE Advanced (2021 Series) - Test 12 (25 May 2021)-img/p1_chem_s1_q02_A.png}{25.5}{0}
\qfigdef{p1_chem_s1_q02_B}{JEE Advanced (2021 Series) - Test 12 (25 May 2021)-img/p1_chem_s1_q02_B.png}{28.0}{0}
\qfigdef{p1_chem_s1_q02_C}{JEE Advanced (2021 Series) - Test 12 (25 May 2021)-img/p1_chem_s1_q02_C.png}{26.5}{0}
\qfigdef{p1_chem_s1_q02_D}{JEE Advanced (2021 Series) - Test 12 (25 May 2021)-img/p1_chem_s1_q02_D.png}{22.3}{0}
\qfigdef{p1_chem_s1_q05_0}{JEE Advanced (2021 Series) - Test 12 (25 May 2021)-img/p1_chem_s1_q05_0.png}{53.4}{0}
\qfigdef{p1_chem_s1_q05_A}{JEE Advanced (2021 Series) - Test 12 (25 May 2021)-img/p1_chem_s1_q05_A.png}{35.5}{0}
\qfigdef{p1_chem_s1_q05_B}{JEE Advanced (2021 Series) - Test 12 (25 May 2021)-img/p1_chem_s1_q05_B.png}{38.2}{0}
\qfigdef{p1_chem_s1_q05_C}{JEE Advanced (2021 Series) - Test 12 (25 May 2021)-img/p1_chem_s1_q05_C.png}{43.0}{0}
\qfigdef{p1_chem_s1_q05_D}{JEE Advanced (2021 Series) - Test 12 (25 May 2021)-img/p1_chem_s1_q05_D.png}{24.5}{0}
\qfigdef{p1_chem_s1_q08_A}{JEE Advanced (2021 Series) - Test 12 (25 May 2021)-img/p1_chem_s1_q08_A.png}{82.7}{0}
\qfigdef{p1_chem_s1_q08_B}{JEE Advanced (2021 Series) - Test 12 (25 May 2021)-img/p1_chem_s1_q08_B.png}{104.6}{0}
\qfigdef{p1_chem_s1_q08_C}{JEE Advanced (2021 Series) - Test 12 (25 May 2021)-img/p1_chem_s1_q08_C.png}{113.9}{0}
\qfigdef{p1_chem_s1_q10_A}{JEE Advanced (2021 Series) - Test 12 (25 May 2021)-img/p1_chem_s1_q10_A.png}{181.0}{0}
\qfigdef{p1_chem_s1_q10_B}{JEE Advanced (2021 Series) - Test 12 (25 May 2021)-img/p1_chem_s1_q10_B.png}{182.5}{0}
\qfigdef{p1_chem_s1_q10_C}{JEE Advanced (2021 Series) - Test 12 (25 May 2021)-img/p1_chem_s1_q10_C.png}{181.8}{0}
\qfigdef{p1_chem_s2_q02_1}{JEE Advanced (2021 Series) - Test 12 (25 May 2021)-img/p1_chem_s2_q02_1.png}{100.7}{0}
\qfigdef{p1_chem_s2_q08_1}{JEE Advanced (2021 Series) - Test 12 (25 May 2021)-img/p1_chem_s2_q08_1.png}{154.8}{0}
\qfigdef{p2_phy_s1_q01_1}{JEE Advanced (2021 Series) - Test 12 (25 May 2021)-img/p2_phy_s1_q01_1.png}{37.3}{0}
\qfigdef{p2_phy_s1_q03_1}{JEE Advanced (2021 Series) - Test 12 (25 May 2021)-img/p2_phy_s1_q03_1.png}{103.2}{0}
\qfigdef{p2_phy_s1_q04_1}{JEE Advanced (2021 Series) - Test 12 (25 May 2021)-img/p2_phy_s1_q04_1.png}{141.8}{0}
\qfigdef{p2_phy_s2_q02_A}{JEE Advanced (2021 Series) - Test 12 (25 May 2021)-img/p2_phy_s2_q02_A.png}{80.6}{0}
\qfigdef{p2_phy_s2_q02_B}{JEE Advanced (2021 Series) - Test 12 (25 May 2021)-img/p2_phy_s2_q02_B.png}{80.3}{0}
\qfigdef{p2_phy_s2_q02_C}{JEE Advanced (2021 Series) - Test 12 (25 May 2021)-img/p2_phy_s2_q02_C.png}{81.1}{0}
\qfigdef{p2_phy_s2_q02_D}{JEE Advanced (2021 Series) - Test 12 (25 May 2021)-img/p2_phy_s2_q02_D.png}{79.8}{0}
\qfigdef{p2_phy_s2_q03_A}{JEE Advanced (2021 Series) - Test 12 (25 May 2021)-img/p2_phy_s2_q03_A.png}{92.0}{0}
\qfigdef{p2_phy_s2_q03_B}{JEE Advanced (2021 Series) - Test 12 (25 May 2021)-img/p2_phy_s2_q03_B.png}{92.1}{0}
\qfigdef{p2_phy_s2_q03_C}{JEE Advanced (2021 Series) - Test 12 (25 May 2021)-img/p2_phy_s2_q03_C.png}{93.3}{0}
\qfigdef{p2_phy_s2_q03_D}{JEE Advanced (2021 Series) - Test 12 (25 May 2021)-img/p2_phy_s2_q03_D.png}{92.6}{0}
\qfigdef{p2_phy_s2_q05_1}{JEE Advanced (2021 Series) - Test 12 (25 May 2021)-img/p2_phy_s2_q05_1.png}{81.5}{0}
\qfigdef{p2_phy_s2_q06_1}{JEE Advanced (2021 Series) - Test 12 (25 May 2021)-img/p2_phy_s2_q06_1.png}{122.7}{0}
\qfigdef{p2_phy_s2_q07_1}{JEE Advanced (2021 Series) - Test 12 (25 May 2021)-img/p2_phy_s2_q07_1.png}{106.3}{0}
\qfigdef{p2_phy_s2_q08_1}{JEE Advanced (2021 Series) - Test 12 (25 May 2021)-img/p2_phy_s2_q08_1.png}{54.1}{0}
\qfigdef{p2_phy_s3_q01_1}{JEE Advanced (2021 Series) - Test 12 (25 May 2021)-img/p2_phy_s3_q01_1.png}{221.3}{0}
\qfigdef{p2_phy_s3_q03_1}{JEE Advanced (2021 Series) - Test 12 (25 May 2021)-img/p2_phy_s3_q03_1.png}{76.1}{0}
\qfigdef{p2_phy_s3_q05_1}{JEE Advanced (2021 Series) - Test 12 (25 May 2021)-img/p2_phy_s3_q05_1.png}{95.0}{0}
\qfigdef{p2_phy_s3_q05_3}{JEE Advanced (2021 Series) - Test 12 (25 May 2021)-img/p2_phy_s3_q05_3.png}{116.5}{0}
\qfigdef{p2_phy_s3_q05_4}{JEE Advanced (2021 Series) - Test 12 (25 May 2021)-img/p2_phy_s3_q05_4.png}{128.8}{0}
\qfigdef{p2_chem_s1_q04_A}{JEE Advanced (2021 Series) - Test 12 (25 May 2021)-img/p2_chem_s1_q04_A.png}{64.2}{0}
\qfigdef{p2_chem_s1_q04_B}{JEE Advanced (2021 Series) - Test 12 (25 May 2021)-img/p2_chem_s1_q04_B.png}{40.6}{0}
\qfigdef{p2_chem_s1_q04_C}{JEE Advanced (2021 Series) - Test 12 (25 May 2021)-img/p2_chem_s1_q04_C.png}{43.9}{0}
\qfigdef{p2_chem_s1_q04_D}{JEE Advanced (2021 Series) - Test 12 (25 May 2021)-img/p2_chem_s1_q04_D.png}{39.4}{0}
\qfigdef{p2_chem_s1_q05_1}{JEE Advanced (2021 Series) - Test 12 (25 May 2021)-img/p2_chem_s1_q05_1.png}{199.3}{0}
\qfigdef{p2_chem_s1_q05_A}{JEE Advanced (2021 Series) - Test 12 (25 May 2021)-img/p2_chem_s1_q05_A.png}{50.4}{0}
\qfigdef{p2_chem_s1_q05_B}{JEE Advanced (2021 Series) - Test 12 (25 May 2021)-img/p2_chem_s1_q05_B.png}{68.3}{0}
\qfigdef{p2_chem_s1_q05_C}{JEE Advanced (2021 Series) - Test 12 (25 May 2021)-img/p2_chem_s1_q05_C.png}{42.6}{0}
\qfigdef{p2_chem_s1_q05_D}{JEE Advanced (2021 Series) - Test 12 (25 May 2021)-img/p2_chem_s1_q05_D.png}{59.9}{0}
\qfigdef{p2_chem_s2_q01_1}{JEE Advanced (2021 Series) - Test 12 (25 May 2021)-img/p2_chem_s2_q01_1.png}{102.3}{0}
\qfigdef{p2_chem_s2_q07_A}{JEE Advanced (2021 Series) - Test 12 (25 May 2021)-img/p2_chem_s2_q07_A.png}{38.4}{0}
\qfigdef{p2_chem_s2_q07_B}{JEE Advanced (2021 Series) - Test 12 (25 May 2021)-img/p2_chem_s2_q07_B.png}{38.0}{0}
\qfigdef{p2_chem_s2_q07_C}{JEE Advanced (2021 Series) - Test 12 (25 May 2021)-img/p2_chem_s2_q07_C.png}{39.7}{0}
\qfigdef{p2_chem_s2_q07_D}{JEE Advanced (2021 Series) - Test 12 (25 May 2021)-img/p2_chem_s2_q07_D.png}{25.7}{0}
\qfigdef{p2_chem_s2_q08_A}{JEE Advanced (2021 Series) - Test 12 (25 May 2021)-img/p2_chem_s2_q08_A.png}{120.7}{0}
\qfigdef{p2_chem_s2_q08_B}{JEE Advanced (2021 Series) - Test 12 (25 May 2021)-img/p2_chem_s2_q08_B.png}{77.4}{0}
\qfigdef{p2_chem_s2_q08_C}{JEE Advanced (2021 Series) - Test 12 (25 May 2021)-img/p2_chem_s2_q08_C.png}{120.7}{0}
\qfigdef{p2_chem_s2_q08_D}{JEE Advanced (2021 Series) - Test 12 (25 May 2021)-img/p2_chem_s2_q08_D.png}{120.4}{0}
\qfigdef{p2_chem_s3_q05_1}{JEE Advanced (2021 Series) - Test 12 (25 May 2021)-img/p2_chem_s3_q05_1.png}{93.0}{0}
\qfigdef{p2_chem_s3_q05_2}{JEE Advanced (2021 Series) - Test 12 (25 May 2021)-img/p2_chem_s3_q05_2.png}{87.2}{0}
\qfigdef{p2_chem_s3_q05_3}{JEE Advanced (2021 Series) - Test 12 (25 May 2021)-img/p2_chem_s3_q05_3.png}{76.2}{0}

\begin{document}
\begin{center}{\Large\bfseries JEE Advanced 12}\end{center}
\begin{multicols}{2}
\papersec{1}
\partsec{1}{PHYSICS}
\sect{SEC--1}
\qn{1} A cubical block of mass $m$ and side length $b$ is placed on a smooth floor. A smooth and rigid rod of length $L$ and with negligible mass is leaning against the block. A sphere of mass $M$ is attached to the upper end of the rod. The lower end of the rod is fixed at point-$O$, but the rod can rotate freely around the point within the vertical plane. Initially the angle between the rod and the floor is $\alpha$ while the system is at rest. After releasing, Choose correct statement(s).
\qfig{p1_phy_s1_q01_1}
\begin{qoptions}
\opt{A}{Speed of the block when angle between rod and the floor is $\beta$ is $b\sqrt{\dfrac{MgL(\sin\alpha-\sin\beta)}{mb^2+ML^2\sin^4\beta}}$}
\opt{B}{Speed of the block when angle between rod and the floor is $\beta$ is $b\sqrt{\dfrac{2MgL(\sin\alpha-\sin\beta)}{mb^2+ML^2\sin^4\beta}}$}
\opt{C}{Angular speed of the rod when angle between rod and the floor is $\beta$ is $\sqrt{\dfrac{2MgL(\sin\alpha-\sin\beta)}{mb^2\mathrm{cosec}^4\beta+ML^2}}$}
\opt{D}{Angular speed of the rod when angle between rod and the floor is $\beta$ is $\sqrt{\dfrac{MgL(\sin\alpha-\sin\beta)}{mb^2+ML^2\cos^4\beta}}$}
\end{qoptions}
\qn{2} Between two parallel long wires $AB$ and $CD$ is a soap film $S$ between another two transverse wires $PQ$ \& $EF$. The separation between $PQ$ \& $EF$ is $x$ which can be varied. The correct statement(s) is/are :
\qfig{p1_phy_s1_q02_1}
\begin{qoptions}
\opt{A}{Force due to surface tension on $PQ$ equals in magnitude to that on $EF$.}
\opt{B}{Force of surface tension on $PQ$ is proportional to the separation $x$.}
\opt{C}{Force of surface tension on $PQ$ is proportional to length of $PQ$.}
\opt{D}{Surface energy depends on $x$.}
\end{qoptions}
\qn{3} In the circuit shown, the switch is closed at $t=0$
\qfig{p1_phy_s1_q03_1}
\begin{qoptions}
\opt{A}{at $t=0$, $I_1=I_2=0$}
\opt{B}{at any time $t=t$, $\dfrac{I_1}{I_2}=\dfrac{L_2}{L_1}$}
\opt{C}{at any time $t=t$, $I_1+I_2=\dfrac\varepsilon R$}
\opt{D}{at $t=\infty$, current flowing through the resistor is $\dfrac\varepsilon R$}
\end{qoptions}
\qn{4} An infinite horizontal square grid of side `$a$' made of uniform material has current $i$ fed into it at node $A$ and this current leaves at $B$. The semi infinite wires used for the entry and exit of current into this grid are shown in the figure-2 which is an enlarged version (not drawn to scale) of the grid in the locality of $AB$ when observed along its plane. The feeding wires are parallel to $AB$ and just above the plane of the grid by a small distance $\delta$ ($\delta\to0$). These feeding wires are thus not visible in figure I. Consider an Amperean circular loop with $AB$ as axis, midpoint of $AB$ as its centre and $R$ as its radius as shown in figure 2. $C$ is the midpoint of square formed by $AB$ as its side shown in figure 1. Take $R=\dfrac a2$
\qfig{p1_phy_s1_q04_1}
\qfig{p1_phy_s1_q04_2}
Magnitude of circulation $\displaystyle\int\vec B\cdot d\vec\ell$ around the shown Amperean loop if $\vec B$ is taken due to grid and feeding wires is $x$.\\
Magnitude of circulation $\displaystyle\int\vec B\cdot d\vec\ell$ around the shown Amperean loop if $\vec B$ is taken of due to entire grid excluding the feeding wires is $y$.\\
Magnitude of $\vec B$ at point $C$ is $z$. Then
\begin{qoptions}
\opt{A}{$x=\dfrac{\mu_0i}2$}
\opt{B}{$y=\mu_0i\left(\dfrac32-\sqrt2\right)$}
\opt{C}{$y=\mu_0i\left(\dfrac1{\sqrt2}-\dfrac12\right)$}
\opt{D}{$z=\dfrac x{\pi a}$}
\end{qoptions}
\qn{5} One mole of an ideal gas undergoes a thermodynamic cyclic process as shown in the figure. The cyclic process consists of an isochoric, an adiabatic process and an isothermal (may not be in the same sequence).
\qfig{p1_phy_s1_q05_1}
\begin{qoptions}
\opt{A}{The ratio of molar specific heat at constant pressure to that at constant volume is $\dfrac{\ln3}{\ln3-\ln2}$.}
\opt{B}{The ratio of molar specific heat at constant pressure to that at constant volume is $\dfrac{\ln2}{\ln3-\ln2}$}
\opt{C}{Magnitude of work done by the gas in the isothermal process is $2P_0V_0\ln(3/2)$.}
\opt{D}{Magnitude of work done by the gas in the isothermal process is $3P_0V_0\ln(3/2)$.}
\end{qoptions}
\qn{6} Two boats are rowing on a pond in the same direction and with the same speed $V$. Each boat sends, through the water (which is at rest), a signal to the other. The frequencies of the generated signals are the same.
\qfig{p1_phy_s1_q06_1}
(i) The time taken by the signals to travel from one boat to the other is different.\\
(ii) The wavelengths received by both the persons in the boats due to the other are different.\\
(iii) The beat frequencies received by both the persons are different.\\
(iv) The beat frequencies received by both the persons are same (which may include zero also)\\
Choose the INCORRECT option(s) :
\begin{qoptions}
\opt{A}{(i), (ii), (iii) are correct}
\opt{B}{(ii), (iii) are correct}
\opt{C}{(iii) is correct}
\opt{D}{(i), (ii), (iv) are correct}
\end{qoptions}
\qn{7} Conducting shell of radius $a$ is surrounded by another spherical conducting shell of radius $b$. The space between them is filled with dielectric material of conductivity $\sigma$ and dielectric constant $K$. The charge $Q_1$ and $Q_2$ are given to the inner and outer shell at time $t=0$ respectively. Charge on outer shell at time $t$ is :
\begin{qoptions}
\opt{A}{$Q_1+Q_2\left(1-\exp\left(-\dfrac{\sigma t}{\varepsilon_0K}\right)\right)$}
\opt{B}{$Q_2-Q_1\left(1-\exp\left(-\dfrac{\sigma t}{\varepsilon_0K}\right)\right)$}
\opt{C}{$Q_1-Q_2\left(1-\exp\left(-\dfrac{\sigma t}{\varepsilon_0K}\right)\right)$}
\opt{D}{$Q_2+Q_1\left(1-\exp\left(-\dfrac{\sigma t}{\varepsilon_0K}\right)\right)$}
\end{qoptions}
\qn{8} A solid sphere is held at rest on the fixed hemispherical bowl made of insulating material such that centre of sphere lies on the horizontal diameter as shown in figure. Total charge $q=5.1\times10^{-1}$ C and mass $m=42$ grams, both are distributed uniformly in the sphere. There exists uniform magnetic field $B$ in space as shown in figure. Consider contact surface to be sufficient rough. Spherical bowl is made of non-magnetic material. Radius of bowl ($R$) $=10$ cm, radius of sphere ($r$) $3$ cm. $g=10$ m/s$^2$. Now if sphere is released
\qfig{p1_phy_s1_q08_1}
\begin{qoptions}
\opt{A}{At an angle $\theta$ with horizontal, velocity of COM of sphere will be $v=\sqrt{\dfrac{10}7g(R-r)\sin\theta}$}
\opt{B}{At an angle $\theta$ with horizontal, velocity of COM of sphere will be $v=\sqrt{\dfrac87g(R-r)\sin\theta}$}
\opt{C}{The limiting value of $B$ (appoximately) so that on releasing sphere, it do not leaves contact with bowl in course of motion will be $2$T.}
\opt{D}{The limiting value of $B$ (appoximately) so that on releasing sphere, it do not leaves contact with bowl in course of motion will be $4$T.}
\end{qoptions}
\qn{9} As shown below, a point light source ($S$) is inside a glass block of refractive index $n_G=\dfrac32$, at a distance $2$ mm from the optical axis and distance $x=3$ cm from the flat surface. A thick plano-convex lens made of transparent plastic material of refractive index $n_P=2$, is placed near the source as shown. The thickness of the lens is $d$, and the radius of convex surface is $R=15$ cm. The width of air gap between the glass block and the lens is `$t$'. The width of air gap $t$ and thickness of lens $d$ is adjusted so that light (paraxial rays) from point source forms always a parallel beam after passing through the lens. Select CORRECT option(s).
\qfig{p1_phy_s1_q09_1}
\begin{qoptions}
\opt{A}{For $d=16$cm, width of air gap $t=5$ cm.}
\opt{B}{For thickness $d=10$cm, width of air gap $t=8$ cm.}
\opt{C}{The maximum adjustable value of $d=26$ cm.}
\opt{D}{The maximum adjustable value of $d=16$ cm.}
\end{qoptions}
\qn{10} A highly thin conducting spherical shell of radius $R_0$ is completely filled with water at $0^\circ$C. Outside surrounding temperature is $-T_0^\circ$C. Latent heat of fusion of water is $L_0$ and thermal conductivity of ice is $K_{ice}$. Assume heat is flowing symmetrically in radial direction. Choose the CORRECT option(s). Density of water is $\rho_0$. Neglect the expansion of water on freezing.
\qfig{p1_phy_s1_q10_1}
\begin{qoptions}
\opt{A}{Time in which complete water will freeze, is given by $\dfrac{\rho_0L_0R_0^2}{6K_{ice}T_0}$}
\opt{B}{Time in which complete water will freeze, is given by $\dfrac{\rho_0L_0R_0^2}{4K_{ice}T_0}$}
\opt{C}{Time in which thickness of ice will become $\dfrac{R_0}2$, is given by $\dfrac{\rho_0L_0R_0^2}{12K_{ice}T_0}$}
\opt{D}{Total heat released by water in freezing, is given by $\dfrac{4\pi\rho_0L_0R_0^3}3$}
\end{qoptions}
\sect{SEC--2}
\qn{1} If the graph of Binding energy ($B/A$) per nucleon versus mass number ($A$) looks like as shown in figure. Fission of nuclei with mass number $150$ into two nuclei of equal mass number will release energy $K\times100$ MeV. find the value of $K$.
\qfig{p1_phy_s2_q01_1}
\qn{2} A man wants to reach directly at point $B$ on the opposite bank of a river flowing at a speed as shown in fig. $V_R$ $=$ Velocity of river with respect to ground.\\
$V_{MR}$ $=$ Velocity of man with respect to river
\qfig{p1_phy_s2_q02_1}
Minimum speed relative to water, the man should have so that he can reach point $B$ will be $\dfrac1{\sqrt k}$ (in m/s). Find the value of $k$.
\qn{3} $AB$ and $CD$ are two uniform resistance wires of lengths $100$ cm and $80$ cm respectively. The connections are shown in the figure. The cell of emf $5$V is ideal while the other cell of emf $E$ has internal resistance $2\,\Omega$ as shown in figure. A length of $20$ cm of wire $CD$ is balanced by $40$ cm of wire $AB$. The emf $E$ in volt is $3x$ then find $x$, if the reading of the ideal ammeter is $2$ A. The other connecting wires have negligible resistance.
\qfig{p1_phy_s2_q03_1}
\qn{4} Snell's law relation for air glass boundary is given by equation $n_{air}\sin\theta_i=n_{glass}\sin\theta_r$. Here angle of refraction in glass is $\theta_r$ and the angle of incident in air is $\theta_i$. The refractive index of air is $n_{air}=1$ (with negligible uncertainty). In an experiment to verify Snell's law, the best fit line on a graph between $\sin\theta_r$ (on the y-axis) and $\sin\theta_i$ (on the x-axis) is given by the equation $y=(6.66\pm0.05)\times10^{-1}x+(0.022\pm0.024)$. The refractive index of the glass, including the error is reported as $1.ab\pm0.0c$, where $a$, $b$ and $c$ are digits after the decimal point. Calculate value of $(a+b-c)$.
\qn{5} A cylindrical beam of light having a variable intensity $I_1$ is incident on the curved surface of a fixed cone of base radius $R$. Intensity of beam changes with radial distance $r$, according to the function $I_1=\dfrac{I_0r^3}{R^3}$. Conical surface is completely reflecting surface. If the force exerted by the light on the cone is of the form $\dfrac{k_1\pi I_0R^2\sin^2\theta}{k_2C}$ where $I_0$ is constant in simplest ratio. Where $k_1$ and $k_2$ are integers and $C$ is the speed of light. Then the value of $(k_1+k_2)$ will be ?
\qfig{p1_phy_s2_q05_1}
\qn{6} Two conducting uncharged kept at a very large distance spheres of radius $R_1$ and $R_2$ ($R_1>R_2$) are connected to a battery with a switch as shown. Light rays of frequency $f$ are incident on the bigger sphere and simultaneously the switch is closed. Work function of bigger sphere is $\phi$. After some time the charge on bigger sphere becomes $q_1$, and on smaller sphere becomes $-q_2$ remains constant there after. Then $\dfrac{q_1}{4\pi\varepsilon_0R_1}+\dfrac{q_2}{K\pi\varepsilon_0R_2}=V$. Find the value of $K$
\qfig{p1_phy_s2_q06_1}
\qn{7} A particle is projected radially outwards from the surface of a planet with an initial speed of $4.0$ km/s. Maximum height attained by the particle is $h=900$ km from the planet's surface. Find the radius $R$ (in km) of the planet taking $g=10$ m/s$^2$ on the surface of the planet. Fill $R/900$ in OMR sheet.
\qn{8} A monochromatic source of light $S$ is placed at focus of a converging lens of focal length $F$. Two identical slits are at `$3F$' distance from the lens and its central line is $\dfrac d2$ distance above optical axis of lens. If $D$ and $d$ have their usual meaning, ($D\gg d\gg\lambda$) is $\dfrac{x\lambda D}{d}$ then fill the value of $x$.
\qfig{p1_phy_s2_q08_1}
\qn{9} An equilateral glass prism kept on a table has refractive index of $\mu=\sqrt2$. It is illuminated by a narrow laser beam having power $P_0$ and wavelength $\lambda$. The path of the laser beam inside the prism is parallel to the base of the prism. The change in reading of weight of the prism due to the incident laser beam is $\dfrac{P_0}c\left(\dfrac{\sqrt x-1}{\sqrt y}\right)$. Find $x+y$ ? $c$ is the speed of light. (Ignore any partial reflection).
\qfig{p1_phy_s2_q09_1}
\qn{10} Two astronauts, $A$ and $B$, both with mass of $60$ kg each, are moving along a straight line in the same direction in a ``weightless'' spaceship. Relative to the spaceship the speed of $A$ is $2$ m/s and that of $B$ is $1$m/s. $A$ is carrying a bag of mass $5$ kg with him. To avoid collision with $B$, $A$ throws the bag with a speed $v$ relative to the spaceship, towards $B$ and $B$ catches it. The minimum value of $v$ is $\left(x+\dfrac{76}{100}\right)$ m/s for which collision is avoided. Find the value of $x$. Assume no change in velocity of spaceship.
\qfig{p1_phy_s2_q10_1}
\partsec{2}{CHEMISTRY}
\sect{SEC--1}
\qn{1} A concentration cell is a galvanic cell in which
\begin{qoptions}
\opt{A}{The electrode material and the solution are made up of same substances but electrode and / or electrolyte concentration are different.}
\opt{B}{Each half cell is made of electrode material and solutions that are composed of different substances and different concentrations.}
\opt{C}{$E^\circ_{\mathrm{Cell}}=0$}
\opt{D}{The nernst equation reduces to $E_{\mathrm{Cell}}=-\dfrac{2.303RT}{nF}\log Q$.}
\end{qoptions}
\qn{2} Which of the following compounds is/are more basic than given structure \qfigi{p1_chem_s1_q01_2} ?
\begin{qoptions}
\opt{A}{\qfigi{p1_chem_s1_q02_A}}
\opt{B}{\qfigi{p1_chem_s1_q02_B}}
\opt{C}{\qfigi{p1_chem_s1_q02_C}}
\opt{D}{\qfigi{p1_chem_s1_q02_D}}
\end{qoptions}
\qn{3} Select the correct statements :
\begin{qoptions}
\opt{A}{Ca$_3$Si$_3$O$_9$ is an example of cyclic silicate}
\opt{B}{Four corner/oxygen atoms per tetrahedron are shared in 3D silicates}
\opt{C}{2D, sheet--like silicates are planar}
\opt{D}{Orthosilicate contains descrete SiO$_4^{4-}$ tetrahedron}
\end{qoptions}
\qn{4} Which of the following is/are colloids ?
\begin{qoptions}
\opt{A}{Smoke}
\opt{B}{Ruby glass}
\opt{C}{Pumice stone}
\opt{D}{Chlorophyll}
\end{qoptions}
\qn{5} When \qfigi{p1_chem_s1_q05_0} is treated with HCl then different carbocations formed can be
\begin{qoptions}
\opt{A}{\qfigi{p1_chem_s1_q05_A}}
\opt{B}{\qfigi{p1_chem_s1_q05_B}}
\opt{C}{\qfigi{p1_chem_s1_q05_C}}
\opt{D}{\qfigi{p1_chem_s1_q05_D}}
\end{qoptions}
\qn{6} Which of the following can be explained on the basis of Fajan's Rules.
\begin{qoptions}
\opt{A}{Ag$_2$S is less soluble than Ag$_2$O.}
\opt{B}{Fe(OH)$_3$ is less soluble than Fe(OH)$_2$.}
\opt{C}{BaCO$_3$ is less soluble than MgCO$_3$.}
\opt{D}{Melting point of AlCl$_3$ is less than that of NaCl.}
\end{qoptions}
\qn{7} Two volatile liquids $A$ \& $B$ are mixed to form an ideal solution. If $P_A^0=100$ mm of Hg \& $P_B^0=150$ mm of Hg then identify the correct option(s) if $10$ moles of $A$ are mixed with $40$ moles of $B$ such that moles vaporised are very small.
\begin{qoptions}
\opt{A}{Total vapour pressure when equilibrium gets established will be $140$ mm of Hg.}
\opt{B}{Mole fraction of $A$ in vapour state will be less than $0.5$.}
\opt{C}{Mole fraction of $B$ in vapour state is equal to $\dfrac67$}
\opt{D}{Mole \% of $B$ in vapour state should be greater than $80\%$.}
\end{qoptions}
\qn{8} Which of the following is/are intermediates or products formed in Perkin reaction ?
\begin{qoptions}
\opt{A}{\qfigi{p1_chem_s1_q08_A}}
\opt{B}{\qfigi{p1_chem_s1_q08_B}}
\opt{C}{\qfigi{p1_chem_s1_q08_C}}
\opt{D}{C$_6$H$_5$---CH $=$ CH --- COOH}
\end{qoptions}
\qn{9} Select the correct IUPAC name for $[(\mathrm{NH_3})_4\mathrm{Co(OH)(NH_2)Co(NH_3)_4}]^{4+}$ :
\begin{qoptions}
\opt{A}{\raggedright $\mu$--Amido-$\mu$-\allowbreak hydroxido\allowbreak octaammine\allowbreak dicobalt(III)ion}
\opt{B}{\raggedright $\mu$--Amido-$\mu$-\allowbreak hydroxido\allowbreak bis\{\allowbreak tetraammine\allowbreak cobalt(III)\}ion}
\opt{C}{\raggedright $\mu$--Amido-$\mu$-\allowbreak hydroxido\allowbreak bis\{\allowbreak tetraammine\allowbreak cobaltate(4+)\}ion}
\opt{D}{\raggedright $\mu$--Amido-$\mu$-\allowbreak hydroxido\allowbreak bis\{\allowbreak tetraammine\allowbreak cobaltate(III)\}ion}
\end{qoptions}
\qn{10} By which of the following reaction sequence can Nylon--5 be prepared.
\begin{qoptions}
\opt{A}{\qfigi[0.9]{p1_chem_s1_q10_A}}
\opt{B}{\qfigi[0.9]{p1_chem_s1_q10_B}}
\opt{C}{\qfigi[0.9]{p1_chem_s1_q10_C}}
\opt{D}{All of these}
\end{qoptions}
\sect{SEC--2}
\qn{1} If the radius of first Bohr orbit of He$^+$ ions is $x$ cm and radius of 4$^{th}$ Bohr orbit of Be$^{3+}$ is $y$ cm, then calculate the value of $y/x$.
\qn{2} \qfig{p1_chem_s2_q02_1}
Give total number of stereoisomers.
\qn{3} Total number of methods which can remove permanent hardness of water.\\
(1) Clark's method\\
(2) Ion-exchange method\\
(3) Synthetic resin method\\
(4) Calgon method\\
(5) Treatment with sodium carbonate
\qn{4} Calculate the pressure exerted (in atm) by $1$ mole of CH$_4$ gas at a temperature of $\dfrac{24}{0.0821}$ K if volume occupied by CH$_4$ molecules is negligible. Given $a=2$ atm lit$^2$ mol$^{-2}$.\\
\textbf{[Give your answer as sum of digits and Fill OMR sheet]}
\qn{5} How many alkenes (C$_5$H$_{10}$) can give 2--methyl--2--butanol on treatment with dil. H$_2$SO$_4$ as major product?
\qn{6} Number of oxides of Nitrogen, which produce HNO$_3$ when dissolved in water.\\
N$_2$O, NO, NO$_2$, N$_2$O$_4$, N$_2$O$_5$.
\qn{7} A metal `M' having atomic mass $31.25$ crystallizes in cubic close packing and it shows ``Schottky defects''. If the edge length of the cubic lattice is $500$ pm and density of the metal is $1.6075$ gm/ml then calculate number of moles of metal atom `M' missing per litre of the crystal.\\
(Given : $1$ amu $=1.67\times10^{-24}$ gm)
\qn{8} \qfig{p1_chem_s2_q08_1}
Give sum of D atom in $A$ and $B$ compound.
\qn{9} What is the oxidation number of metal in binary compound used as catalyst for conversion of SO$_2$ to SO$_3$ in contact process for the making of H$_2$SO$_4$ ?
\qn{10} Among the following, total number of possible metals present in anode mud, which is obtained by electrolytic refining of copper.\\
Fe, Ag, Au, Pt, Pb, Zn
\partsec{3}{MATHEMATICS}
\sect{SEC--1}
\qn{1} If $\vec a,\vec b,\vec c$ are unit vectors such that $\vec a\cdot\vec b=0$ and $(\vec a-\vec c)\cdot(\vec b+\vec c)=0$ \& $\vec c=\lambda\vec a+\mu\vec b+\omega(\vec a\times\vec b)$ for some scalars $\lambda,\mu,\omega$. Then
\begin{qoptions}
\opt{A}{$\mu=\lambda+1$}
\opt{B}{$\lambda^2+\mu^2+\omega^2=2$}
\opt{C}{$\lambda^2+\mu^2+\omega^2=1$}
\opt{D}{$\lambda=\mu+1$}
\end{qoptions}
\qn{2} If `$\alpha$' is minimum value of $f(x)=\dfrac{x^4-2x^2+3}{(x^2+1)^2}$ and `$\beta$' is/are value(s) of $x$ at which $f(x)$ attains its minimum, then
\begin{qoptions}
\opt{A}{$\alpha$ is rational number}
\opt{B}{$\beta$ is/are irrational number(s)}
\opt{C}{$\alpha=\dfrac13$}
\opt{D}{$|\beta|=\sqrt2$}
\end{qoptions}
\qn{3} If $A(z_1)$ is a variable point in argand plane such that $z_1\bar z_1=5$, $0<\arg(z_1)<\pi$. Also, $B(z_2)$ \& $C(z_3)$ are two fixed points in argand plane satisfying $z^4-6z^2+25=0$ \& $-\pi<\arg z<0$.
\begin{qoptions}
\opt{A}{Internal angle bisector of $\angle A$ of $\Delta ABC$ will always pass through a fixed point.}
\opt{B}{Circumcentre of $\Delta ABC$ is origin}
\opt{C}{Orthocentre $(z_4)$ of $\Delta ABC$ satisfies $|z_4+2i|=5$. $(i=\sqrt{-1})$}
\opt{D}{Orthocentre $(z_4)$ of $\Delta ABC$ satisfies $|z_4-2i|=\sqrt5$. $(i=\sqrt{-1})$}
\end{qoptions}
\qn{4} A biased coin with probability of getting heads double that of probability of getting tails is tossed. If this coin is tossed repeatedly until we get two successive tails \& $P_n$ denotes the probability that this experiment is completed on $n^{\text{th}}$ toss. Then :
\begin{qoptions}
\opt{A}{$P_2=\dfrac19$}
\opt{B}{$P_2=\dfrac29$}
\opt{C}{$P_4=\dfrac5{81}$}
\opt{D}{$P_4=\dfrac2{27}$}
\end{qoptions}
\qn{5} A line $y=2x+c$ intersects the parabola $y^2=4x$ in two distinct points $A$ \& $B$, such that normals at $A$ and $B$ intersect on the principal axis of parabola. Then
\begin{qoptions}
\opt{A}{$AB$ is perpendicular to principal axis.}
\opt{B}{Exactly one of $A$ \& $B$ is vertex of parabola}
\opt{C}{$AB$ subtends right angle at focus}
\opt{D}{Foot of perpendicular from focus upon $AB$ divides $AB$ in the ratio $1:2$ or $2:1$}
\end{qoptions}
\qn{6} Consider a plane $P:2x+y-2z=6$ and line $L:\dfrac x2=\dfrac{y-\lambda}{-1}=\dfrac z2$. Plane \& line intersect at a point $A$ on $xy$--plane. Then :
\begin{qoptions}
\opt{A}{Acute angle between line \& plane is $\sin^{-1}\!\left(\dfrac19\right)$}
\opt{B}{Line is perpendicular to the plane}
\opt{C}{Co-ordinates of $A$ are $(0,6,0)$}
\opt{D}{Coordinates of $A$ are $(6,6,0)$}
\end{qoptions}
\qn{7} If $I_n=\displaystyle\int_{-\pi}^{\pi}\dfrac{1}{1+2^{\sin\left(\frac x2\right)}}\left(\dfrac{\sin\left(\frac{nx}2\right)}{\sin\left(\frac x2\right)}\right)^2dx$ for $n\in W$. Which of the following statements are true ?
\begin{qoptions}
\opt{A}{$I_n=I_{n+1}$ for all $n>1$}
\opt{B}{$I_0,I_1,I_2,\ldots,I_n$ forms an AP.}
\opt{C}{$\displaystyle\sum_{n=0}^{9}I_{2n}=90\pi$}
\opt{D}{$\displaystyle\sum_{m=0}^{10}I_m=55\pi$}
\end{qoptions}
\qn{8} If the equation $(x^2+x+2)^2-(a-3)(x^2+x+2)(x^2+x+1)+(a-4)(x^2+x+1)^2=0$ has at least one real root. Then `$a$' can be ?
\begin{qoptions}
\opt{A}{6}
\opt{B}{$\dfrac{11}2$}
\opt{C}{5}
\opt{D}{7}
\end{qoptions}
\qn{9} A circle cuts the parabola $x^2=16y$ at $O(0,0)$, $A(-4,1)$, $B(x_1,y_1)$ and $C(x_2,y_2)$. If normals at $A$ and $B$ to the parabola intersect at $C$, then-
\begin{qoptions}
\opt{A}{$x_1+y_2=49$}
\opt{B}{$x_2+y_1=100$}
\opt{C}{$(0,-8)$, $A$ \& $B$ are collinear}
\opt{D}{$(0,8)$, $A$ \& $C$ are collinear}
\end{qoptions}
\qn{10} In a game a person $P$ at point $(x,y)$ can jump to any one of the following points relatively: $(x+1,y)$, $(x+2,y)$, $(x,y+1)$ or $(x,y+2)$. If $P$ is at $(0,0)$, then he can reach
\begin{qoptions}
\opt{A}{$(2,2)$ in 14 ways}
\opt{B}{$(2,3)$ in 32 ways}
\opt{C}{$(2,2)$ in 9 ways}
\opt{D}{$(2,3)$ in 26 ways}
\end{qoptions}
\sect{SEC--2}
\qn{1} Let $f'(x)=-\dfrac{f(x)}{\sqrt{f^2(x)+g^2(x)}}$ \& $g'(x)=1-\dfrac{g(x)}{\sqrt{f^2(x)+g^2(x)}}$ $\forall\,x\in R$, $g(0)=0$, $f(0)=10$ and $\displaystyle\lim_{x\to\infty}f(x)=0$, then value of $\displaystyle\lim_{x\to\infty}g(x)$ is
\qn{2} Consider two points $A(1,1,2)$ and $B(3,5,4)$. If locus of point $P$ on $xy$--plane such that $|PA-PB|$ is minimum, is represented by set $\{(x,y,0),\ x+ky=\ell\}$, then the value of $|k-\ell|$ is
\qn{3} Let $I_n=\displaystyle\int_0^{\pi/2}\tan^n\!\left(\dfrac x2\right)dx$; $n\in N$ then the value of $\displaystyle\lim_{n\to\infty}n(I_n+I_{n+2})$ is
\qn{4} Consider a $3\times3$ matrix $A=\{a_{ij}\}$ where $a_{ij}=\dfrac1{i+j-1}$ $\forall\,i,j\in\{1,2,3\}$. If matrix $B=\{b_{ij}\}$ of order $3\times3$ exists such that $A^{-1}=3B$, then the value of $\dfrac{|B|}{a_{23}\cdot b_{22}}$ is (where $|X|$ denotes determinant value of matrix $X$)
\qn{5} 4 married couples are to be seated on a round table such that no two men are seated together. If $P=$ Number of ways no wife is seated adjacent to their husbands and $Q=$ Number of ways no wife is seated to the right seat of their husbands, then value of $\left[\dfrac QP\right]$ is (where $[.]$ denotes greatest integer function)
\qn{6} Number of ordered pairs $(x,y)$ with $0\le x\le\pi$ \& $\pi\le y\le2\pi$ satisfying the equation $(\sin x+\sqrt3\cos x)(\sin(x+y)+\sin(x-y))=4$ is
\qn{7} Given that $\alpha,\beta,a,b$, are in AP. $\alpha,\beta,c,d$ are in GP and $\alpha,\beta,e,f$ are in HP. If $b,d,f$ are in GP then find $\dfrac{2(\alpha^6-\beta^6)}{\alpha\beta(\alpha^4-\beta^4)}$, $0<\alpha<\beta$.
\qn{8} If sum of the series ${}^{12}C_5+7\cdot{}^{12}C_6+21\cdot{}^{12}C_7+35\cdot{}^{12}C_8+35\cdot{}^{12}C_9+21\cdot{}^{12}C_{10}+7\cdot{}^{12}C_{11}+{}^{12}C_{12}$ is equal to ${}^{2n+1}C_{12}$, then $n$ is equal to
\qn{9} Consider a triangle $ABC$ with $A(0,0)$, $B(2\sqrt3,0)$ and $C(\sqrt3,3)$. If $L$ is locus of the points from which tangents of equal lengths can be drawn to its incircle \& excircle opposite to vertex $A$, then the minimum distance of $L$ from origin is
\qn{10} If $m,n$ are the roots of quadratic equation $ax^2+bx+c=0$, then roots of the equation $(4a+2b+c)x^2+(4a-b-2c)x+(a-b+c)=0$ can be expressed as $\dfrac{m+e}{m+f}$, $\dfrac{n+g}{n+h}$, $m\ne2$, $n\ne2$, then the value of $efgh$, is
\papersec{2}
\partsec{1}{PHYSICS}
\sect{SEC--1}
\qn{1} The pulley of radius $r$ as shown in figure has a moment of inertia $I$ about its axis and mass $m$. The time period of vertical oscillation of its centre of mass is, if the spring has spring constant $k$ and the string does not slip over the pulley, is :-
\qfig{p2_phy_s1_q01_1}
\begin{qoptions}
\opt{A}{$2\pi\sqrt{\dfrac{\left(\frac{I}{r^2}+m\right)}{(2k)}}$}
\opt{B}{$2\pi\sqrt{\dfrac{\left(\frac{I}{r^2}-m\right)}{(4k)}}$}
\opt{C}{$2\pi\sqrt{\dfrac{\left(\frac{I}{r^2}+m\right)}{(4k)}}$}
\opt{D}{$2\pi\sqrt{\dfrac{\left(\frac{I}{r^2}-m\right)}{(2k)}}$}
\end{qoptions}
\qn{2} A string of $0.1$ kg/m and length $0.6$ m is fixed at both ends and stretched such that it has a tension of 40 N. The string vibrates in three segments with maximum amplitude of 0.5 cm. Then maximum transverse velocity of a particle which is at a distance $\dfrac1{30}$ m from fixed end is :
\begin{qoptions}
\opt{A}{$25\pi$ cm/s}
\opt{B}{25 cm/s}
\opt{C}{$50\pi$ cm/s}
\opt{D}{50 cm/s}
\end{qoptions}
\qn{3} In the given figure two identical waves each of intensity $I_0$; 1 and 2 are superimposed. The resulting intensity is :
\qfig{p2_phy_s1_q03_1}
\begin{qoptions}
\opt{A}{$I_0$}
\opt{B}{$2\sqrt2I_0$}
\opt{C}{$4I_0$}
\opt{D}{$2I_0$}
\end{qoptions}
\qn{4} Plane sound waves of wavelength $0.12$ m are incident on two narrow slits in a box with non reflecting walls, as shown. At a distance of $5.0$ m from the center of the slits, a first-order maximum occurs at point $P$, which is $3.0$ m from the central maximum. The distance between the slits is most nearly :-
\qfig{p2_phy_s1_q04_1}
\begin{qoptions}
\opt{A}{0.07 m}
\opt{B}{0.09 m}
\opt{C}{0.16 m}
\opt{D}{0.20 m}
\end{qoptions}
\qn{5} Two thin symmetrical lenses of different nature (one converging and other diverging) have equal radii of curvature $R=20$ cm. The lenses are put close together and immersed in water. The converging focal length of the system is 24 cm. Find the difference between refractive indices of the two lenses. (refractive index for water $=\dfrac43$)
\begin{qoptions}
\opt{A}{$(\mu_1-\mu_2)=\dfrac1{24}$}
\opt{B}{$(\mu_1-\mu_2)=\dfrac1{20}$}
\opt{C}{$(\mu_1-\mu_2)=\dfrac59$}
\opt{D}{$(\mu_1-\mu_2)=\dfrac43$}
\end{qoptions}
\sect{SEC--2}
\qn{1} Steel blocks $A$ and $B$, which have equal masses, are at $(T_A)_i=300^\circ$C and $(T_B)_i=400^\circ$C. Block $C$, with $m_C=2m_A$ is at $(T_C)_i=350^\circ$C. Blocks $A$ and $B$ are placed in contact, allowed to come to equilibrium and isolated. Then they are placed in contact with block $C$. After equilibrium is reached, final temperatures of $A$, $B$ and $C$ is $T_A$, $T_B$ and $T_C$, then
\begin{qoptions}
\opt{A}{$T_A=\sqrt{T_BT_C}$}
\opt{B}{$T_A=\dfrac{T_BT_C}{T_B+T_C}$}
\opt{C}{$\dfrac{T_A+T_B}2=T_C$}
\opt{D}{$2T_A-T_B=T_C$}
\end{qoptions}
\qn{2} A particle initially at rest on a smooth horizontal surface, is acted on by a constant horizontal force at time $t=0$. Then if $W=$ work done, $t=$ time and $v=$ speed of the particle, the nearly correct graph(s) is/are :-
\begin{qoptions}
\opt{A}{\qfigi{p2_phy_s2_q02_A}}
\opt{B}{\qfigi{p2_phy_s2_q02_B}}
\opt{C}{\qfigi{p2_phy_s2_q02_C}}
\opt{D}{\qfigi{p2_phy_s2_q02_D}}
\end{qoptions}
\qn{3} Two identical balls (1 and 2) moving with equal speed collide on a frictionless surface. Collision may or may not be elastic (may or may not be head on). Which of the following diagrams is/are a physically possible situation ?
\begin{qoptions}
\opt{A}{\qfigi{p2_phy_s2_q03_A}}
\opt{B}{\qfigi{p2_phy_s2_q03_B}}
\opt{C}{\qfigi{p2_phy_s2_q03_C}}
\opt{D}{\qfigi{p2_phy_s2_q03_D}}
\end{qoptions}
\qn{4} The spherically symmetric electric field is given by $\vec E=\begin{cases}\alpha\vec r & 0\le r<R\\ 2\alpha\vec r & r>R\end{cases}$. Where $\alpha$ is a positive constant. Mark the correct option(s) :
\begin{qoptions}
\opt{A}{This cannot be induced electric field due to change in magnetic field.}
\opt{B}{This electric field can be generated by superposition of a continuously varying finite charge distribution and uniform surface charge distribution.}
\opt{C}{The graph of electric potential is continuous and have a kink.}
\opt{D}{The outside region (i.e. $r>R$) also contain some charge.}
\end{qoptions}
\qn{5} A gas takes part in three thermal processes in which it is heated from the same initial state to the same final temperature. The process are shown on the $P$-$V$ diagram by straight lines 1-2, 1-3 and 1-4. $Q_1$, $Q_2$ and $Q_3$ are amount of heat supplied to the gas in the process $1\to2$, $1\to3$, $1\to4$, respectively.
\qfig{p2_phy_s2_q05_1}
\begin{qoptions}
\opt{A}{$Q_3>Q_2$}
\opt{B}{$Q_1<Q_2$}
\opt{C}{$Q_3>Q_1$}
\opt{D}{$Q_3>Q_2>Q_1$}
\end{qoptions}
\qn{6} A massless frame is present in uniform magnetic field and a block of mass $m$ hangs on the frame as shown in figure. When a constant current $I>\dfrac{mg}{\ell B}$ is maintained in the frame, it gets displaced by `$h$' in some time interval. If $\varepsilon$ is the emf of battery then which of the following is/are correct ?
\qfig{p2_phy_s2_q06_1}
\begin{qoptions}
\opt{A}{Work done by magnetic force is $I\ell Bh$.}
\opt{B}{Extra Work done by battery due to present of magnetic field is $I\ell Bh$.}
\opt{C}{Velocity of block when it gets displacement by $h$ is $\sqrt{\dfrac{2(I\ell B-mg)h}{m}}$.}
\opt{D}{Net Force on charges in the segment $PQ$ while the frame moves up is in the vertical direction.}
\end{qoptions}
\qn{7} Curved surface of a vessel has shape of a truncated cone having semi vertex angle $37^\circ$. Vessel is full of water (density $\rho=1000$ kg/m$^3$) upto a height of 13 cm and is placed on a smooth horizontal plane. Upper surface is opened to atmosphere. A hole of $1.5$ cm$^2$ is made on curved wall at a height of 8 cm from bottom as shown in figure. Area of water surface in the vessel is large as compared to the area of hole.
\qfig{p2_phy_s2_q07_1}
\begin{qoptions}
\opt{A}{Initial velocity of efflux is 1 m/sec}
\opt{B}{Initial horizontal range of water jet from point $B$ is $6.65$ cm}
\opt{C}{Horizontal force required to keep the vessel in static equilibrium is $0.15$ N}
\opt{D}{Horizontal force required to keep the vessel in static equilibrium is $0.12$ N.}
\end{qoptions}
\qn{8} For the biprism experiment shown in the figure, the fringe width increases when ($S$ is a slit from which monochromatic waves are incident on the biprism).
\qfig{p2_phy_s2_q08_1}
\begin{qoptions}
\opt{A}{biprism is moved towards the slit.}
\opt{B}{entire apparatus is submerged in a liquid having R.I. less than that of prism.}
\opt{C}{a biprism having smaller angle $\alpha$ is used.}
\opt{D}{the slit width is reduced.}
\end{qoptions}
\sect{SEC--3}
\qn{1} Block $A$ of mass 4 kg is placed on the rough surface of a long block $B$ of mass 6 kg. Block $B$ is placed on smooth horizontal surface. Blocks are given the velocities as shown in figure.
\qfig{p2_phy_s3_q01_1}
\end{multicols}
\begin{center}\small
\begin{tabular}{|p{1cm}|p{8cm}|p{1cm}|p{2cm}|}\hline
 & \textbf{List-I} & & \textbf{List-II}\\\hline
(P) & Common velocity (in m/s) of blocks when sliding between $A$ and $B$ stops. & (1) & 108\\\hline
(Q) & Velocity (in m/s) of centre of mass at shown moment. & (2) & 80\\\hline
(R) & Magnitude of total work done (in joule) by friction on $B$. & (3) & 8\\\hline
\end{tabular}
\end{center}
\noindent\qn{2} A particle of mass $m$ is projected with velocity $v$ and at an angle $\theta$ with horizontal then match the following list.\par\nopagebreak\smallskip
\begin{center}\small
\begin{tabular}{|p{1cm}|p{8cm}|p{1cm}|p{3cm}|}\hline
 & \textbf{List-I} & & \textbf{List-II}\\\hline
(P) & Magnitude of Change in momentum between point of projection and point at which it strikes on horizontal surface & (1) & $mv\cos\theta$\\\hline
(Q) & Radius of curvature at highest point & (2) & $2mv\sin\theta$\\\hline
(R) & Minimum magnitude of momentum is & (3) & $\dfrac{u^2\cos^2\theta}{g}$\\\hline
\end{tabular}
\end{center}
\begin{multicols}{2}
\qn{3} A monotomic ideal gas is taken through a cyclic process $ABCDA$. $ABCDA$ is circular path. If work done is denoted by $W$ then match the List-I with List-II.
\qfig{p2_phy_s3_q03_1}
\end{multicols}
\begin{center}\small
\begin{tabular}{|p{1cm}|p{3cm}|p{1cm}|p{3cm}|}\hline
 & \textbf{List-I} & & \textbf{List-II}\\\hline
(P) & $AB$ & (1) & $W>0$\\\hline
(Q) & $AC$ & (2) & $W<0$\\\hline
(R) & $CD$ & (3) & $W=0$\\\hline
\end{tabular}
\end{center}
\noindent\qn{4} Match the following dimensions of list-I to list-II\par\nopagebreak\smallskip
\begin{center}\small
\begin{tabular}{|p{1cm}|p{4cm}|p{1cm}|p{4cm}|}\hline
 & \textbf{List-I} & & \textbf{List-II}\\\hline
(P) & Thermal resistance & (1) & $[MT^{-3}\theta^{-4}]$\\\hline
(Q) & Stefan's constant & (2) & $[M^{-1}L^{-2}T^{3}\theta]$\\\hline
(R) & Wien's constant & (3) & $[L\theta]$\\\hline
\end{tabular}
\end{center}
\newpage\noindent\qn{5}\par\nopagebreak\smallskip
\begin{center}\small
\begin{tabular}{|p{1cm}|p{6.5cm}|p{1cm}|p{6cm}|}\hline
 & \textbf{List-I} & & \textbf{List-II}\\\hline
(P) & \qfigi{p2_phy_s3_q05_1}\rule{0pt}{70pt} & (1) & Time constant is the time in which current drops to 37\% of its initial value.\\\hline
(Q) & \qfigi{p2_phy_s3_q05_3}\rule{0pt}{70pt} & (2) & Time constant is the time in which current increases to 63\% of its steady state value.\\\hline
(R) & \qfigi{p2_phy_s3_q05_4}\rule{0pt}{80pt} & (3) & Time constant is the time in which current increases to 37\% of its steady state value.\\\hline
\end{tabular}
\end{center}
\begin{multicols}{2}
\partsec{2}{CHEMISTRY}
\sect{SEC--1}
\qn{1} 200 ml of 0.005M AgNO$_3$ reacts with 300 ml of 0.01M KCl. If $K_{sp}$ of AgCl is $1.8\times10^{-10}$, then maximum concentration of Ag$^+$ in the mixture is -
\begin{qoptions}
\opt{A}{$6\times10^{-3}$M}
\opt{B}{$4\times10^{-3}$M}
\opt{C}{$2\times10^{-3}$M}
\opt{D}{$4.5\times10^{-8}$M}
\end{qoptions}
\qn{2} For particles having same kinetic energy, the de-broglie wavelength is -
\begin{qoptions}
\opt{A}{Directly proportional to mass}
\opt{B}{Inversely proportional to speed}
\opt{C}{Inversely proportional to square root of mass}
\opt{D}{Inversely proportional to square root of speed}
\end{qoptions}
\qn{3} Which gas produces brown coloured solution when it is passed through aqueous solution of FeSO$_4$?
\begin{qoptions}
\opt{A}{NO}
\opt{B}{CO}
\opt{C}{NH$_3$}
\opt{D}{H$_2$S}
\end{qoptions}
\qn{4} Which of the following compound give cerric ammonium nitrate test ?
\begin{qoptions}
\opt{A}{\qfigi{p2_chem_s1_q04_A}}
\opt{B}{\qfigi{p2_chem_s1_q04_B}}
\opt{C}{\qfigi{p2_chem_s1_q04_C}}
\opt{D}{\qfigi{p2_chem_s1_q04_D}}
\end{qoptions}
\qn{5}
\qfig{p2_chem_s1_q05_1}
\begin{qoptions}
\opt{A}{\qfigi{p2_chem_s1_q05_A}}
\opt{B}{\qfigi{p2_chem_s1_q05_B}}
\opt{C}{\qfigi{p2_chem_s1_q05_C}}
\opt{D}{\qfigi{p2_chem_s1_q05_D}}
\end{qoptions}
\sect{SEC--2}
\qn{1} Consider the following graph for the reaction $A\to P$ :
\qfig{p2_chem_s2_q01_1}
choose the correct statement(s)
\begin{qoptions}
\opt{A}{The given reaction is a second order reaction}
\opt{B}{The given reaction is a third order reaction}
\opt{C}{Initial half life of reaction is 4 min}
\opt{D}{Initial concentration of A is 2M}
\end{qoptions}
\qn{2} Which of the following statements is/are true for critical temperature ?
\begin{qoptions}
\opt{A}{It is the highest temperature at which liquid and vapour coexist.}
\opt{B}{Above critical temperature, there is no distinction between liquid and gas phases}
\opt{C}{At critical point, the density of the sytem is zero}
\opt{D}{At critical point, the viscosity of the system is zero.}
\end{qoptions}
\qn{3} In which of the following metal extraction, following steps are involved?\\
(i) Roasting \quad (ii) Carbon reduction \quad (iii) Electrolytic refining
\begin{qoptions}
\opt{A}{Lead from galena}
\opt{B}{Magnesium from carnelite}
\opt{C}{Iron from haematite}
\opt{D}{Zinc from sphalerite}
\end{qoptions}
\qn{4} Select CORRECT statement(s).
\begin{qoptions}
\opt{A}{O$_3$ and O$^{2-}$ both are diamagnetic}
\opt{B}{Out of O$_2$, O$_2^+$ and O$_2^-$, O$_2$ have maximum magnetic moment (spin only)}
\opt{C}{Out of O$_2$, O$_2^+$ and O$_3$, O$_2^+$ have least O--O bond length}
\opt{D}{Out of O$_2$, O$_2^{2+}$ and O$_2^{2-}$, only O$_2$ is paramagnetic.}
\end{qoptions}
\qn{5} Which of the following produce both Cr$_2$O$_3$ and O$_2$ in product side?
\begin{qoptions}
\opt{A}{Reaction of Al$_2$O$_3$ with Cr}
\opt{B}{Heating CrO$_3$}
\opt{C}{Heating (NH$_4$)$_2$Cr$_2$O$_7$}
\opt{D}{Heating K$_2$Cr$_2$O$_7$}
\end{qoptions}
\qn{6} Which of the following is/are reducing sugar
\begin{qoptions}
\opt{A}{Sucrose}
\opt{B}{Aldohexose}
\opt{C}{Aldotriose}
\opt{D}{Ketotriose}
\end{qoptions}
\qn{7} Which of the following compound react with HNO$_2$ ?
\begin{qoptions}
\opt{A}{\qfigi{p2_chem_s2_q07_A}}
\opt{B}{\qfigi{p2_chem_s2_q07_B}}
\opt{C}{\qfigi{p2_chem_s2_q07_C}}
\opt{D}{\qfigi{p2_chem_s2_q07_D}}
\end{qoptions}
\qn{8} Which of the following reaction is/are gives primary alcohol as a major product ?
\begin{qoptions}
\opt{A}{\qfigi{p2_chem_s2_q08_A}}
\opt{B}{\qfigi{p2_chem_s2_q08_B}}
\opt{C}{\qfigi{p2_chem_s2_q08_C}}
\opt{D}{\qfigi{p2_chem_s2_q08_D}}
\end{qoptions}
\sect{SEC--3}
\qn{1}
\end{multicols}
\begin{center}\small
\begin{tabular}{|p{1cm}|p{7cm}|p{1cm}|p{4cm}|}\hline
 & \textbf{List-I} (Reaction) & & \textbf{List-II} ($\Delta_rH$ is known as)\\\hline
(P) & C(diamond) $\to$ C(g) & (1) & $\Delta_{\text{lattice}}H$\\\hline
(Q) & CH$_3$COOH(aq.) + OH$^-$(aq.) $\to$ CH$_3$COO$^-$(aq.) + H$_2$O(l) & (2) & $\Delta_{\text{atomisation}}H$\\\hline
(R) & NaCl(s) $\to$ Na$^+$(g) + Cl$^-$(g) & (3) & $\Delta_{\text{neutralisation}}H$\\\hline
\end{tabular}
\end{center}
\noindent\qn{2}\par\nopagebreak\smallskip
\begin{center}\small
\begin{tabular}{|p{1cm}|p{5cm}|p{1cm}|p{4cm}|}\hline
 & \textbf{List-I} & & \textbf{List-II}\\\hline
(P) & H$_2$(g) + I$_2$(g) $\rightleftharpoons$ 2HI(g) & (1) & $K_P=K_C(RT)^1$\\\hline
(Q) & PCl$_5$(g) $\rightleftharpoons$ PCl$_3$(g) + Cl$_2$(g) & (2) & $K_C=K_P(RT)^1$\\\hline
(R) & 2NO$_2$(g) $\rightleftharpoons$ N$_2$O$_4$(g) & (3) & $K_p=K_C$\\\hline
\end{tabular}
\end{center}
\noindent\qn{3}\par\nopagebreak\smallskip
\begin{center}\small
\begin{tabular}{|p{1cm}|p{5cm}|p{1cm}|p{3cm}|}\hline
 & \textbf{List-I} & & \textbf{List-II}\\\hline
(P) & Complex follow Sidgwick rule & (1) & [Pt(en)ClBr]\\\hline
(Q) & Complex have geometrical isomer & (2) & Na[Co(CO)$_4$]\\\hline
(R) & Complex have optical isomer & (3) & [Cr(en)$_2$Cl$_2$]Cl\\\hline
\end{tabular}
\end{center}
\noindent\qn{4}\par\nopagebreak\smallskip
\begin{center}\small
\begin{tabular}{|p{1cm}|p{3cm}|p{1cm}|p{7cm}|}\hline
 & \textbf{List-I} & & \textbf{List-II}\\\hline
(P) & ZnSO$_4$ solution & (1) & Gives chocolate brown ppt with K$_4$[Fe(CN)$_6$] solution\\\hline
(Q) & CuSO$_4$ solution & (2) & Gives white ppt with dil HCl\\\hline
(R) & AgNO$_3$ solution & (3) & Gives Reinmann's green with Co(NO$_3$)$_2$ solution followed by heating\\\hline
\end{tabular}
\end{center}
\noindent\qn{5}\par\nopagebreak\smallskip
\begin{center}\small
\begin{tabular}{|p{1cm}|p{6cm}|p{1cm}|p{3cm}|}\hline
 & \textbf{List-I}\newline Reaction with their major product & & \textbf{List-II}\newline Missing reagent\\\hline
(P) & \qfigi{p2_chem_s3_q05_1}\rule{0pt}{26pt} & (1) & NaOH / CaO\\\hline
(Q) & \qfigi{p2_chem_s3_q05_2}\rule{0pt}{36pt} & (2) & Na / Dry ether\\\hline
(R) & \qfigi{p2_chem_s3_q05_3}\rule{0pt}{60pt} & (3) & Cl$_2$ / h$\nu$\\\hline
\end{tabular}
\end{center}
\begin{multicols}{2}
Missing reagent for reactions $PQR$ respectively?
\partsec{3}{MATHEMATICS}
\sect{SEC--1}
\qn{1} If $f_1(x)=e^x$ and $f_{n+1}(x)=e^{f_n(x)}$ $\forall\,n\in N$. Then $\dfrac{d}{dx}\left(f_n(x)\right)$ is equal to-
\begin{qoptions}
\opt{A}{$f_n(x)$}
\opt{B}{$f_n(x)f_1(x)$}
\opt{C}{$f_n(x)f_{n-1}(x)\ldots f_1(x)$}
\opt{D}{$f_n(x)f_{n-1}(x)\ldots f_1(x)e^x$}
\end{qoptions}
\qn{2} A normal is drawn at point $P$ of an ellipse, which meets the major axis at point $G$. If $SP=2$ and $S'P=8$, then the ratio of areas of $\Delta PSG$ and $\Delta PS'G$ is (where $S$ and $S'$ are foci)
\begin{qoptions}
\opt{A}{$4:1$}
\opt{B}{$1:1$}
\opt{C}{$1:2$}
\opt{D}{$1:4$}
\end{qoptions}
\qn{3} At the beginning of the year 2019, India's population was 130 crore and it is increasing at constant rate of 1\% annually. Right now the mobile numbers are of 10 digits starting with digits 6,7,8 or 9 only, as permitted by TRAI (a regulating body). Find the maximum number of years the current pattern of mobile numbers can run, assuming two mobile numbers per member of population (use $\log_{(1.01)}\left(\dfrac{20}{13}\right)=43.29$)
\begin{qoptions}
\opt{A}{44 years}
\opt{B}{43 years}
\opt{C}{22 years}
\opt{D}{21 years}
\end{qoptions}
\qn{4} Two distinct function $f(x)$ and $g(x)$ are chosen randomly from a set of functions $\left\{x,\,x^2,\,\sin\left(\dfrac1x\right),\,\cos\left(\dfrac1x\right)\right\}$. Then the probability that $\displaystyle\lim_{x\to0}f(x)\,g(x)$ exist, is
\begin{qoptions}
\opt{A}{$\dfrac56$}
\opt{B}{$\dfrac1{12}$}
\opt{C}{$\dfrac13$}
\opt{D}{$\dfrac16$}
\end{qoptions}
\qn{5} If $|z|=1$, then locus of $w=\dfrac{z-1}{z+1}\ (z\ne-1)$ is (where $z$ and $w$ are complex numbers)
\begin{qoptions}
\opt{A}{a circle}
\opt{B}{imaginary axis}
\opt{C}{real axis}
\opt{D}{a ray}
\end{qoptions}
\sect{SEC--2}
\qn{1} Choose the correct statement(s)-
\begin{qoptions}
\opt{A}{The largest integer in the range of function $(\sin^{-1}x)^2-(\cos^{-1}x)^2$ is 2}
\opt{B}{The smallest integer in the range of function $(\sin^{-1}x)^2-(\cos^{-1}x)^2$ is $-7$}
\opt{C}{The value of `$x$' satisfying $\sin^{-1}x+2\cos^{-1}(-x)=2\pi$, is greater than $\dfrac12$}
\opt{D}{The value of `$x$' satisfying $\sin^{-1}x+2\cos^{-1}(-x)=2\pi$, is less than $\dfrac12$}
\end{qoptions}
\qn{2} If $A$ is a square matrix such that $A^2=I$, then $\det\left((A-I)^3+(A+I)^3-7A\right)$ can be- (where $\det(A)$ denotes determinant of square matrix $A$)
\begin{qoptions}
\opt{A}{0}
\opt{B}{$-1$}
\opt{C}{1}
\opt{D}{2}
\end{qoptions}
\qn{3} A particle moves over a path, whose position vector at any time $t$, is given by $\vec r(t)=|t|\,\hat i+|\sin\pi t|\,\hat j+\sqrt[3]{t-2}\,\hat k$. A path $\vec r(t)$ is called differentiable at time $t$ if its all components are differentiable with respect to `$t$' at that point of time. So $\vec r(t)$ is differentiable at $t=$
\begin{qoptions}
\opt{A}{3}
\opt{B}{1}
\opt{C}{2}
\opt{D}{3.5}
\end{qoptions}
\qn{4} If $f(x)=\begin{cases}2\sqrt x & ;\ x\in[0,1)\\ p & ;\ x=1\\ -\ell nx & ;\ x>1\end{cases}$. Then at $x=1$, $f(x)$ has local
\begin{qoptions}
\opt{A}{maxima if $p=2$}
\opt{B}{maxima if $p=1$}
\opt{C}{minima if $p=0$}
\opt{D}{minima if $p=-1$}
\end{qoptions}
\qn{5} If $\displaystyle\int e^x(x+1)^2dx=(x^2+px+q)e^x+c$ (where $c$ is constant of integration and $p,q\in R$), then-
\begin{qoptions}
\opt{A}{$pq=0$}
\opt{B}{$p+q=1$}
\opt{C}{$q=0$}
\opt{D}{$p=0$}
\end{qoptions}
\qn{6} If {\small\[\int_0^\infty\dfrac{x^2dx}{(x^2+a^2)(x^2+b^2)(x^2+c^2)}=\dfrac{\pi}{2(a+b)(b+c)(c+a)}\]}
$a,b,c\in R$ further,
{\small\begin{gather*}I=\int_0^\infty\dfrac{dx}{(x^2+4)(x^2+9)}\quad\text{and}\\ J=\int_0^\infty\dfrac{dx}{(x^2+1)(x^2+4)(x^2+9)},\end{gather*}}
then
\begin{qoptions}
\opt{A}{$I=\dfrac\pi{60}$}
\opt{B}{$I=\dfrac\pi{30}$}
\opt{C}{$J=\dfrac\pi{120}$}
\opt{D}{$J=\dfrac\pi{240}$}
\end{qoptions}
\qn{7} Consider the curve\\
$S=(\sin\theta)x^2+2(\cos\theta)xy+(\sin\theta)y^2+2(\cos\theta)x+2(\cos\theta)y+\sin\theta=0$\\
Now, choose the correct statement(s)
\begin{qoptions}
\opt{A}{If $\theta=\dfrac{3\pi}4$, then $S=0$ represents a pair of straight lines.}
\opt{B}{If $S=0$ represents a pair of real straight lines, then $\theta=n\pi+\dfrac\pi4$, $n\in Z$}
\opt{C}{If $S=0$ represents a pair of real straight lines, then they are parallel.}
\opt{D}{There exist some $\theta$ for which $S=0$ represents no real locus.}
\end{qoptions}
\qn{8} Consider a family of curves in which tangent at any point always passes through $(1,0)$. Then
\begin{qoptions}
\opt{A}{Differential equation governing the given family is $\dfrac{dy}{dx}=\dfrac{x-1}y$}
\opt{B}{The general equation of given family is $y=m(x-1)$ ; $m$ is arbitrary constant}
\opt{C}{A curve of the given family passing through $(2,1)$ is $y=x-1$}
\opt{D}{$x$-axis belongs to the given family.}
\end{qoptions}
\sect{SEC--3}
\qn{1}
\end{multicols}
\begin{center}\small
\begin{tabular}{|p{1cm}|p{8.5cm}|p{1cm}|p{1.5cm}|}\hline
 & \textbf{List-I} & & \textbf{List-II}\\\hline
(P) & If $x,y>0$ such that $2x+y=3$, then maximum value of $x^2y$ is & (1) & 1\\\hline
(Q) & The value of $[-2\sin25^\circ]^2$ is equal to \newline (where $[.]$ denotes greatest integer function) & (2) & 2\\\hline
(R) & The value of $\left|\dfrac{\sqrt2+\sqrt3\,w+\sqrt5\,w^2}{\sqrt5+\sqrt2\,w+\sqrt3\,w^2}\right|$ is equal to \newline ($w$ is imaginary cube root of unity) & (3) & 4\\\hline
\end{tabular}
\end{center}
\noindent\qn{2} List-I contains vector equation in unknown $\vec r$ and List-II contains number of solutions i.e. number of such possible $\vec r$.\par\nopagebreak\smallskip
\begin{center}\small
\begin{tabular}{|p{1cm}|p{6cm}|p{1cm}|p{3cm}|}\hline
 & \textbf{List-I} & & \textbf{List-II}\\\hline
(P) & $\vec r\times\hat i=\hat i$ & (1) & 0\\\hline
(Q) & $\vec r\cdot(\hat i+\hat j+\hat k)=1$ and $|\vec r|=\dfrac1{\sqrt3}$ & (2) & 1\\\hline
(R) & $\left|\vec r-\hat i\right|=1$ & (3) & more than 1\\\hline
\end{tabular}
\end{center}
\noindent\qn{3}\par\nopagebreak\smallskip
\begin{center}\small
\begin{tabular}{|p{1cm}|p{8cm}|p{1cm}|p{1.5cm}|}\hline
 & \textbf{List-I} & & \textbf{List-II}\\\hline
(P) & The number of roots of equation $\sin^2x-5\sin x+6<0$ ; $x\in[0,2\pi]$, is & (1) & 0\\\hline
(Q) & The eccentricity of a hyperbola $24x^2-y^2=24$, is & (2) & 3\\\hline
(R) & The angle between the tangents drawn from $(1,4)$ to the parabola $y^2=4x$ is $\dfrac{2\pi}k$, where $k$ is equal to & (3) & 5\\\hline
\end{tabular}
\end{center}
\noindent\qn{4}\par\nopagebreak\smallskip
\begin{center}\small
\begin{tabular}{|p{1cm}|p{8.5cm}|p{1cm}|p{1.5cm}|}\hline
 & \textbf{List-I} & & \textbf{List-II}\\\hline
(P) & If the equation of a plane containing the co-planar lines $\dfrac{x-2}1=\dfrac{y-9}2=\dfrac{z-a}3$ and $\dfrac{x+3}1=\dfrac{y-1}{-2}=\dfrac{z+2}3$ $(a\in R)$ is $3x-z+k=0$, then $k$ is equal to & (1) & 7\\\hline
(Q) & $AB$ is a chord of the circle $x^2+y^2=9$. The tangent at $A$ and $B$ intersect at $C$. If $M(1,2)$ is the mid-point of $AB$, then the area of $\Delta ABC$ is $\dfrac p{\sqrt5}$, where $p$ is equal to & (2) & 8\\\hline
(R) & If number of words that can be formed using all the letters of the word `ALLEN' is $N$, then sum of digits of $N$ is & (3) & 6\\\hline
\end{tabular}
\end{center}
\noindent\qn{5}\par\nopagebreak\smallskip
\begin{center}\small
\begin{tabular}{|p{1cm}|p{8.5cm}|p{1cm}|p{1.5cm}|}\hline
 & \textbf{List-I} & & \textbf{List-II}\\\hline
(P) & If $(1+x+x^2)^{10}=\displaystyle\sum_{r=0}^{20}a_rx^r$ then $a_5-{}^{10}C_1a_4+{}^{10}C_2a_3-{}^{10}C_3a_2+{}^{10}C_4a_1-{}^{10}C_5a_0$ is equal to & (1) & 0\\\hline
(Q) & The period of $f(x)=\dfrac{e^x}{e^{[x]}}$, is (where $[.]$ denotes greatest integer function) & (2) & 5\\\hline
(R) & The maximum slope of tangent to the curve $y=-x^3+3x^2+2x$ is & (3) & 1\\\hline
\end{tabular}
\end{center}
\begin{multicols}{2}
\end{multicols}
\end{document}
